\chapter{Mengentheoretische Folgerungen}
\label{b07guide:ch:images}

In diesem Kapitel sei stets $f\colon A\twoheadrightarrow B$.
Es seien $V,W\subseteq B$.

\section{Zielteilmengen durch Urbilder erkennen}

\begin{theorem}[Gleichheit, Inklusion und echte Inklusion]
\label{b07guide:thm:reflection}
Es gelten
\[
 \begin{aligned}
 f^{-1}[V]=f^{-1}[W]&\Longleftrightarrow V=W,\\
 f^{-1}[V]\subseteq f^{-1}[W]&\Longleftrightarrow V\subseteq W,\\
 f^{-1}[V]\subsetneq f^{-1}[W]&\Longleftrightarrow V\subsetneq W.
 \end{aligned}
\]
Hier bezeichnet $\subsetneq$ echte Teilmengeninklusion.
\end{theorem}

\begin{theorem}[Leere Mengen, volle Mengen und Disjunktheit]
\label{b07guide:thm:tests}
Es gelten
\[
 \begin{aligned}
 f^{-1}[V]=\varnothing&\Longleftrightarrow V=\varnothing,\\
 f^{-1}[V]=A&\Longleftrightarrow V=B,\\
 f^{-1}[V]\cap f^{-1}[W]=\varnothing
       &\Longleftrightarrow V\cap W=\varnothing.
 \end{aligned}
\]
Für jede gegebene Familie $(V_i)_{i\in I}$ von Teilmengen von $B$
sind die $V_i$ genau dann paarweise disjunkt, wenn die Urbilder
$f^{-1}[V_i]$ paarweise disjunkt sind.
\end{theorem}

\begin{theorem}[Überdeckungen und Komplemente erkennen]
\label{b07guide:thm:cover}
Für jede gegebene Familie $(V_i)_{i\in I}$ von Teilmengen von $B$ gilt
\[
 \bigcup_{i\in I}V_i=B
 \quad\Longleftrightarrow\quad
 \bigcup_{i\in I}f^{-1}[V_i]=A.
\]
Ferner gilt
\[
 W=B\setminus V
 \quad\Longleftrightarrow\quad
 f^{-1}[W]=A\setminus f^{-1}[V].
\]
Beide Aussagen schließen leere Mengen und leere Familien ein.
\end{theorem}

\begin{definition}[Durchschnitt relativ zu einer Grundmenge]
Für $E_i\subseteq X$, $i\in I$, sei
\[
 \bigcap_{i\in I}^{X}E_i
   =\{x\in X\mid\forall i\in I\ (x\in E_i)\}.
\]
Der leere Durchschnitt ist dabei $X$.
\end{definition}

\begin{theorem}[Gemeinsame Zielwerte durch gemeinsame Urbilder erkennen]
\label{b07guide:thm:intersection}
Für jede gegebene Familie $(V_i)_{i\in I}$ von Teilmengen von $B$ gilt
\[
 f\left[\bigcap_{i\in I}^{A}f^{-1}[V_i]\right]
     =\bigcap_{i\in I}^{B}V_i.
\]
Insbesondere gilt
\[
 \bigcap_{i\in I}^{A}f^{-1}[V_i]\ne\varnothing
 \quad\Longleftrightarrow\quad
 \bigcap_{i\in I}^{B}V_i\ne\varnothing.
\]
Für $I=\varnothing$ lautet die erste Gleichheit $f[A]=B$.
Die Aussage behauptet keine Schnittregel für Bilder
beliebiger Teilmengen von $A$.
\end{theorem}

\section{Aussagen auf die Ausgangsmenge übertragen}

\begin{theorem}[Übertragung beschränkter Quantoren]
\label{b07guide:thm:quantifiers}
Für jede mengentheoretische Formel $\varphi(b)$ mit festen
Mengenparametern gelten
\[
 \begin{aligned}
 \forall b\in B\ \varphi(b)
   &\Longleftrightarrow \forall a\in A\ \varphi(f(a)),\\
 \exists b\in B\ \varphi(b)
   &\Longleftrightarrow \exists a\in A\ \varphi(f(a)).
 \end{aligned}
\]
Für eine Formel $\psi(b,c)$ und eine beliebige Menge $C$ gilt außerdem
\[
 \begin{aligned}
 \forall b\in B\;\exists c\in C\ \psi(b,c)
 &\Longleftrightarrow
   \forall a\in A\;\exists c\in C\ \psi(f(a),c),\\
 \exists b\in B\;\forall c\in C\ \psi(b,c)
 &\Longleftrightarrow
   \exists a\in A\;\forall c\in C\ \psi(f(a),c).
 \end{aligned}
\]
Hier wird keine Funktion behauptet, die für alle $b$
gleichzeitig einen Zeugen $c$ auswählt.
\end{theorem}